% 07-discussion.tex 中文对照版（译自 sections/07-discussion.tex；不加 \label）
\section{讨论}

\paragraph{为什么是控制面，而不是更好的记忆。} 在不变量 I1 之下，更锐利的记忆也是更锐利的危险：检索到的陈述越忠实，就越诱人把它当作当前事实呈现。把设计精力从检索质量转向准入与降级，直接针对失效面，代价是让检索质量大体停留在原来的位置。

\paragraph{张力 1：多数 fail-open 与唯一 fail-closed 点。} 这种不对称是刻意的——安全失败必须阻断，功能失败必须保持可用——但它有代价：功能失败可以是静默的。补偿手段是可观测性，而不是对称性。每个被丢弃的块都携带原因码，降级被记录而不是被抛出异常，每一轮都发出可离线回放的事件。这并不能让静默失败变得不可能，但能让它在事后可归因。

\paragraph{张力 2：冻结的评测仪器与学习闭环。} 如果评测仪器可写，学习闭环学到的第一件事就是改考卷。冻结指标定义、回放集与护栏阈值阻止了这一点，但也压缩了学习能起作用的可用空间。我们自己的搜索说明了代价：在 972 组参数配置中，没有任何配置把目标函数改善到超过实际显著性阈值（$\Delta < 5\times10^{-4}$），而六个一级检索参数都变得可辨识。我们把它作为关于评测仪器的方法学结果报告，而不是关于性能的结果。

\paragraph{张力 3：自包含插件与单一事实来源。} 宿主插件必须能在没有 Python 运行时的条件下安装，因此控制参数存在于三处：权威规则文件、宿主 patch 配置与插件默认值。66 键比对可以检出漂移，但检出漂移并不等于证明行为等价。对这次移植的诚实描述是“带文档化差异清单的重新实现”，而不是“绑定”。

\paragraph{代价。} 当前工件中有四项可见代价。第一，预算：注入的块与对话历史竞争，而这一权衡由本身就是 guess（猜测值）的静态权重支配。第二，延迟：控制面被设计为停留在自己设定的小预算内，但那个预算只是目标，不是测量值。第三，维护：同一套语义的两份实现必须靠手工保持一致。第四，审计体量：八轮真实对话产生 308 条审计事件，这是不可忽视的存储与检查负担。

\paragraph{可迁移性。} 应当区分两个层次。机制层——阈值语义、失败语义、冻结仪器的纪律——与宿主无关，可以独立于任何 API 来规定。接线层——基于 section 的 prompt 注入、钩子名称、审批 seam——则特定于某一个宿主。因此移植需要重新推导接线，并逐条对照新宿主复核已记录的差异。

\paragraph{可达性与效应量。} 本文的证据确立三件事：机制在真实宿主中可达、契约可执行、失败可审计且可回放。它并未在大规模真实流量上确立显著的正向效应量。这一空缺是结构性的，而不是疏忽：测量这样的效应需要我们不能制造的流量，以及我们不应在别人的部署上执行的干预。这恰恰是这些工件开源的原因——引擎、回放集、混沌语料、预注册协议与分析工具——以便有条件获得更多访问权的人能够复现或反驳这一测量。我们宁愿公开仪器与实测到的零结果，也不愿给出一个我们无法辩护的数字。

\paragraph{方法学教训。} 把“我们不能声称什么”当作一等工件，成本很低且反复有用：一份假设台账、给每个数字打上 provenance 标签、以及一件冻结的评测仪器，使过度声称成为机制性的失败，而不是修辞上的失败。我们把这一做法作为独立于本系统的建议推荐。
